\section{New industries, fads, and waves}\label{sec:waves}

This section asks when and where the leading penalty of
Section~\ref{sec:gate} lived: across registration cohorts, across the
technology classes, across the themes each era's leading filings were
drawn from, in the filings that name the internet, and inside surging
themes. The outcome throughout is the same one --- failure at the
five-year proof of continued use --- and every sample is settled cohorts
(registrations whose ninth anniversary is inside the record), so nothing
here is censored.

\subsection{The penalty weakened in the middle of the sample and was largest at the start}\label{sec:eras}

Scores exist from 1995, which makes registration cohorts from 1996
evaluable, and their outcomes at the five-year proof have been settled for
two decades. Figure~\ref{fig:cohortslopes} draws, for each cohort, the
change in survival for a one-standard-deviation increase in lead and in
atypicality, the two entered jointly and standardized within Nice class
and cohort. Pooled across classes, the lead coefficient is significantly
negative in every settled cohort except 2002--2005 --- $-1.9$ to $-2.3$pp
per standard deviation for the 1996--2000 cohorts, $+0.4$ to $-0.2$
across 2002--2005, and $-0.5$ to $-1.4$ from 2006 on. In the software
class the middle of the sample is not a pause but a reversal: the
coefficient runs from $-6.8$pp for the 1996 cohort to significantly
\emph{positive} ($+0.9$ to $+1.2$) in 2002, 2003 and 2005, then back
through $-4.0$ in 2012. The atypicality coefficient never changes sign
anywhere: more atypical registrations survive more in every cohort of
both panels, by $+0.6$ to $+2.5$pp per standard deviation pooled and
$+2.9$ to $+5.7$ in software.

\begin{figure}[h]\centering
\includegraphics[width=\textwidth]{results/fig_cohort_slopes.png}
\caption{Survival of the five-year proof of continued use, per
registration cohort, as continuous coefficients: the change in survival,
in percentage points, for a one-standard-deviation increase in lead (red)
and in atypicality (blue), entered jointly, both standardized within Nice
class and cohort. Left: all classes pooled ($n = 3{,}150{,}442$ settled
registrations); right: software and electronics alone. Bands are 95\%
intervals on OLS standard errors, uncorrected for the within-owner
correlation of outcomes (0.38), so they understate the uncertainty
somewhat.}
\label{fig:cohortslopes}
\end{figure}

Stated as cohort contrasts, the swing runs from $+6.9$pp for the 1996
cohort through $+10.1$pp for 1999 --- roughly twice the largest cohort
contrast observed since --- to the collapse, and back to $+1.8$ to
$+5.6$pp for every cohort from 2008. Neither prosecution-speed selection
in the earliest cohorts nor the cohort indexing drives it: holding
pendency in a common band leaves the series unchanged, and indexed by
filing year the contrast is $+7.9$ to $+9.6$pp for 1995--1999 filings
against $-0.95$ for filings made in 2000 (supplementary material).

Splitting the four classes that carry most technology filings --- 009,
035, 038 and 042 --- against the other forty-one separates a large moving
component from a small steady one (Figure~\ref{fig:eras}):

\begin{center}\small
\begin{tabular}{lrr}
\toprule
Filing years & Technology classes & All others \\
\midrule
1995--1999 & $+14.18$ (0.36) & $+4.48$ (0.27) \\
2000--2004 & $-1.37$ (0.35) & $+1.21$ (0.24) \\
2005--2007 & $+2.04$ (0.43) & $+1.99$ (0.28) \\
2008--2014 & $+9.02$ (0.25) & $+1.32$ (0.17) \\
\bottomrule
\end{tabular}
\end{center}

\noindent Outside technology the penalty is a flat $+1$ to $+4$pp in
every era, including the one in which the pooled series reads as zero.
Inside technology it swings by fifteen points: strongly positive for
filings made into the late 1990s, negative for those made in the five
years after, and positive again from 2008. The pooled quiet of 2002--2005
is a technology reversal partly offsetting a non-technology positive, not
an absence.

What this design cannot say is which cycle the swing belongs to. Marks
filed 1995--1999 reach their five-year proof in 2001--2007 and marks
filed 2000--2004 reach theirs in 2006--2013, so a story in which unusual
language is punished because it was written during an investment boom and
a story in which it is punished because the proof falls due during a bust
both fit these numbers.

\begin{figure}[h]\centering
\includegraphics[width=0.86\textwidth]{results/gate_era_profile.png}
\caption{The leading penalty at the five-year proof by registration
cohort, technology classes (009, 035, 038, 042) against all others. Each
point is the failure rate of the most leading fifth of that cohort's
registrations minus that of the most lagging fifth, lead quintiles cut
within cohort and group, no controls, with 95\% intervals, on settled
registrations. Outside technology the penalty is small and almost flat
across two decades; inside technology it swings from $+16$pp for the 2000
cohort to $-3$pp for 2003 and back to $+12$pp for 2012. The dotted line
marks the 2002 cut that the cohort samples of Section~\ref{sec:gate}
use.}
\label{fig:eras}
\end{figure}

\subsection{Three of the four classes reverse, and the leading language changes theme}\label{sec:techswing}

The four-class split pools industries that do not move together. Taking
them one at a time, and then assigning quintiles inside progressively
finer cells (Table~\ref{tab:techswing}):

\begin{table}[h]\centering\small
\caption{The leading penalty at the five-year proof inside technology, by
class and by cell. Top-minus-bottom lead-quintile contrast in failure
(pp, SE in parentheses), settled registrations, by filing era. The upper
block assigns quintiles within each class and era. The lower block pools
the four classes and assigns quintiles within the era (raw), within
class$\times$registration-cohort cells, and within
theme$\times$class$\times$registration-year cells, where a filing's theme
is the one its description loads most heavily on under the production
$T = 50$ model. $n = 1{,}248{,}591$.}
\label{tab:techswing}
\input{results/gate_era_tech_classes.tex}
\end{table}

\noindent Electrical and software goods (009) and telecommunications
(038) follow the pooled pattern. Business services (035), the class that
holds online retail, reverses hardest: $-9.7$pp for filings made
2000--2004 against $+4.6$pp in the five years before. Scientific and
computer services (042) never reverses. Its contrast is $+6.7$pp in the
era in which the other three are negative and between $+8$ and $+17$pp in
every other, and by registration cohort it never falls below $+3$pp
(Figure~\ref{fig:techswing}). The reversal is an event in goods and in
business services; the class that sells technology as a service sat it
out.

Assigning quintiles within class and registration year rather than within
the era pool changes little. Assigning them within theme, class and
registration year changes a great deal: the late-1990s contrast falls
from $+14.5$ to $+4.8$pp, the post-2008 contrast from $+9.2$ to $+4.5$pp,
and the 2000--2004 reversal from $-3.6$ to $-1.3$pp. Two thirds of the
first and half of the second are between themes rather than within them.
The leading fifth of each era was drawn from particular themes, and those
themes carried their own failure rates; the penalty for being leading
\emph{inside} a theme is a steadier $+2.5$ to $+5$pp, the same order as
the penalty outside technology.

Which themes is a matter of record (Table~\ref{tab:techthemes}; the
table's last row is the all-themes pooled contrast within
theme$\times$class$\times$cohort). In 1995--1999 the computer, software,
information and data services theme held 38\% of technology filings and
75\% of the leading fifth, and failed at 64.9\% --- that one theme is
most of the late-1990s penalty. In 2000--2004 the leading language had
moved to electronic media while the internet-services wording became
lagging, and inside several themes the lagging fifth now failed more than
the leading one. From 2008 the leading fifth over-represents downloadable
and online software four to one --- the application store's vocabulary,
the highest-failure theme (67.2\%) with the largest within-theme contrast
($+8.4$pp pooled across cohorts). The full theme tour is in the
supplementary material.

\begin{table}[h]\centering\footnotesize
\caption{The leading penalty at the five-year proof by theme within the
four technology classes. Themes with at least 20{,}000 settled
registrations, labelled by their five highest-probability terms;
quintiles assigned within theme and era. Cells with fewer than 3{,}000
registrations are blank. The final row pools all themes, with quintiles
cut within theme$\times$class$\times$cohort.}
\label{tab:techthemes}
\input{results/gate_era_tech_themes.tex}
\end{table}

\begin{figure}[h]\centering
\includegraphics[width=0.92\textwidth]{results/gate_era_tech_themes.png}
\caption{The leading penalty at the five-year proof by registration
cohort, one line per technology class. Each point is the failure rate of
the most leading fifth of that class and cohort's registrations minus
that of the most lagging fifth --- lead quintiles cut within class and
cohort, no controls --- with 95\% intervals, on settled registrations.
Three classes cross zero between 2001 and 2006; scientific and computer
services does not.}
\label{fig:techswing}
\end{figure}

So the swing has two layers. Across the whole period, a leading filing is
a few points more likely to fail than a lagging one in the same theme,
class and cohort, in technology as elsewhere. On top of that, each era's
leading vocabulary was concentrated in one or two themes whose marks
failed at their own rate: internet services in the 1990s, application
software after 2008. In the trough between, the language that had been
leading became lagging, and the marks still written in it failed with it.
The design-based estimates of Table~\ref{tab:gate_lpm} hold class and
cohort fixed but not theme, so they carry both layers; a reader who wants
only the within-theme penalty should take the lower block of
Table~\ref{tab:techswing} as the guide to its size.

\subsection{Internet as a theme: web-based filings against their class}\label{sec:internet}

A filing is internet-bearing when its goods/services text names the
internet, the web, a website, online or on-line provision, or electronic
commerce. The rule reads the text rather than the class, so a web
retailer filing in apparel is caught and a semiconductor maker filing in
class 009 is not. Internet-bearing filings are 16.9\% of scored
registrations: 58\% in telecommunications, 39--41\% in advertising and
retail, scientific and technical services, and legal and personal
services, under 5\% in most goods classes.
Figure~\ref{fig:webscatter} puts every class on two axes at once --- how
much more often its internet-bearing filings register, and how much more
often they fail the five-year proof; the supplementary material
tabulates the underlying rates for every class.

\begin{figure}[h]\centering
\includegraphics[width=0.8\textwidth]{results/fig_internet_scatter.png}
\caption{Each Nice class's internet-bearing filings against the rest of
the class, on both steps at once. Horizontal: registration rate of
internet-bearing class-records minus the rest of the class, filings
1995--2018. Vertical: failure at the five-year proof of continued use,
internet-bearing minus the rest, registrations 2002--2018. Classes with
fewer than 500 internet-bearing registrations are dropped. Classes in the
upper-right quadrant grant internet filings more often and lose them
more often.}
\label{fig:webscatter}
\end{figure}

Three patterns are in the figure. In the industrial goods classes,
internet-bearing registrations fail the five-year proof far more often
than their class: machinery 50.9\% against 38.7\%, chemicals 46.1\%
against 39.1\%, vehicles 49.6\% against 45.1\%. A web-based business
registering in an engineering class fails at the rate of a software
registration (51.4\%), not at the rate of its class. In the consumer
goods classes the sign reverses --- internet-bearing registrations in
clothing fail 52.4\% against 58.7\% for the class, in household utensils
44.5\% against 50.5\% --- so online sellers of consumer goods outlast the
class's other registrants. In every service class the internet-bearing
filings fail more often, by one to eleven points, most in legal and
personal services (60.1\% against 48.7\%). At the registration step the
industrial-goods classes sit to the right of zero: internet-bearing
applications there reach registration 59.5\% of the time against 56.0\%
for the rest. For those classes the two steps move in opposite
directions --- more likely to be granted, then more likely to be
cancelled for non-use six years later --- which is the pattern
Section~\ref{sec:gate_conclusion} predicts when the examiner reads
distance from class language and the market reads whether the product
stayed in commerce.

The failure gap has a common shape in event time. Dating each class's
internet arrival at the first filing year the pattern reaches 2\% of its
filings, the gap starts near $+7$ points in the first observable cohorts
and settles between $+1$ and $+4$ nine to eleven years after arrival,
shrinking in six of the eight classes observable in both phases; calendar
time disguises this, since machinery's arrival was 2008 and its peak
2012--2015 (supplementary material). The class's clock starts when the
vocabulary arrives, not when the internet did --- the shape falling entry
costs followed by normalization would leave.

The internet split also locates the era swing of
Section~\ref{sec:eras}. Rerunning the per-cohort lead coefficients of
Figure~\ref{fig:cohortslopes} separately for internet-bearing and other
registrations (Figure~\ref{fig:webslopes}): the mid-sample reversal
belongs to the internet-bearing filings, whose leading registrations
survived \emph{more} in the 2001--2005 cohorts ($+1.3$ to $+2.5$pp per
standard deviation) while non-internet technology registrations sat near
zero. The other two episodes are not the internet's in this textual
sense. The late-1990s penalty is present in technology registrations that
never name the web ($-3.4$ to $-4.8$pp for the 1996--2000 cohorts) ---
consistent with Table~\ref{tab:techthemes}, where that era's deadly
vocabulary is computer and software services, words that predate explicit
internet naming. And the post-2008 penalty sits in both splits at once
($-2$ to $-3$pp each). Outside technology, the non-internet series is
the steady $-0.4$ to $-1.4$pp throughout. What the internet changed, on
this cut, is the middle: for half a decade after the crash, the web
filings whose language led their class were the ones that lasted.

\begin{figure}[h]\centering
\includegraphics[width=\textwidth]{results/fig_cohort_slopes_internet.png}
\caption{The per-cohort lead coefficient of Figure~\ref{fig:cohortslopes},
split by the curated internet text pattern: change in survival of the
five-year proof per standard deviation of lead, atypicality held,
standardized within Nice class and cohort. Left: the four technology
classes; right: all other classes. Bands are 95\% intervals on OLS
standard errors; cohort cells under 2{,}000 registrations are dropped,
which is why the internet-bearing lines start late.}
\label{fig:webslopes}
\end{figure}

The Nice Classification has twice faced this question and twice declined
to make the internet a class.\footnote{The eighth edition (2002) split
the old class 42 into classes 42--45 by kind of service; the twelfth
(2023) placed downloadable virtual goods in class 9 and distributed
metaverse services across the existing service classes.} The survival
pattern bears on the choice: internet-bearing filings in a goods class
fail at rates closer to the technology classes than to their own, so the
technology overlay carries more information for this outcome than the
class assignment --- and Nice classes used as industry controls, as this
paper and the literature use them, do not partition the internet economy.

\subsection{Surges: joining a wave}\label{sec:surge}

Genuinely new themes are rare (Section~\ref{sec:three_scorings}); the
observable event is a surge, an established theme's share doubling, over
its own three prior years, to a level that matters. The one corpus-scale
episode --- online software doubling to 2.4\% of all filings in 1999 ---
failed the five-year proof $8.6$ points above class and year, $8.5$ net
of the filings' own lead and atypicality, with no gradient in arrival
order inside the wave ($-0.2$pp, SE $1.9$). Class-scale surges are common
(962 episodes at a 1\% within-class threshold) and cost little pooled:
$+0.65$pp within class and year, $+1.27$ net of lead and atypicality,
with episodes differing in sign among themselves. Within waves, entry
order does not matter; larger waves are not deadlier; and growing waves
carry more excess failure than receding ones ($+4.3$ against $+1.0$
points, suggestive on 61 episodes) --- a pattern consistent with
selection at the proof rather than crowding out. The surge figure, the
episode table, and the wave-level facts are in the supplementary
material.

\subsection{What the waves show}\label{sec:waves_conclusion}

The leading penalty is steady and small outside technology and episodic
inside it, and most of the episode is between themes: each era's leading
vocabulary was concentrated in one or two themes that failed at their own
rate, while the within-theme penalty holds at $+2.5$ to $+5$pp
throughout. Web-based filings fail hardest in each class's own first
internet years and converge toward their class within a decade. Joining a
surging theme costs little on average and nothing for arriving early
within the wave; the one corpus-scale wave, online software in 1999, is
the exception at every margin. What a fad costs, on this record, is
mostly a matter of which theme the fad was, not of when the filer joined
it.
